%%% General package setup %%%

% Languages %
% Multilingual typesetting for LuaLaTeX (replacing Babel)
% (see https://ctan.org/pkg/polyglossia)
\usepackage{polyglossia}

% Set default language to American English and register
% Dutch and French as other languages
\setdefaultlanguage[variant=usmax,ordinalmonthday=false]{english}
\setotherlanguage{dutch}
\setotherlanguage{french}

% Context-sensitive quotations
% (see https://ctan.org/pkg/csquotes)
\usepackage[
autostyle=true, % Automatically adapt style to current language
]{csquotes}


% Order-dependent, early %
% Mathtools (includes amsmath, extending it)
% (see https://ctan.org/pkg/mathtools)
\usepackage[
intlimits, % Integral limits below/above in display math (passed to amsmath)
]{mathtools}

% Inline math italics correction and always show manual tags
\mathtoolsset{mathic=true,showmanualtags=true}

% Protrude equation numbers into the margin
% (See https://tony-zorman.com/posts/phd-typesetting.html#the-margin)
\makeatletter
\let\oldmaketag@@@\maketag@@@%
\def\oldtagform@#1{\oldmaketag@@@{(\ignorespaces#1\unskip\@@italiccorr)}}
\renewcommand{\eqref}[1]{\textup{\oldtagform@{\ref{#1}}}}
\newlength{\width@@}
\def\maketag@@@#1{%
  \hbox{%
    \hskip1sp\m@th\llap{%
      \normalfont%
      \normalcolor%
      #1%
      \settowidth{\width@@}{#1}%
      \Ifthispageodd{\hspace{-\the\width@@-\the\marginparsep}}{\hspace{\textwidth+\the\marginparsep}}%
    }%
  }%
}%

% Remove minimum separation length between tag and equation within text/display
% area, as the tag fully protrudes into the margin anyways
% This aims to prevent the tag from being shifted downward when the displayed
% equation overlaps what would be the *original* position of the tag
% (without the above margin protrusion)
\renewcommand{\mintagsep}{0pt}
\makeatother


% Fonts and (Micro-)Typography %
% Fontspec for font selection in LuaLaTeX
% (see https://ctan.org/pkg/fontspec)
\usepackage{fontspec}

% Libertinus (OTF) fonts, similar to Linux Libertine (Mono) and Biolinum,
% but with added math font.
% (see https://ctan.org/pkg/libertinus-otf)
\usepackage[
osf,p, % Old-style, proportional numerals (in text)
sb, % Semi-bold version
]{libertinus-otf}
% Useful Libertinus features:
% - \Lctosc, \LCtoSC+, \LCtoSC+, for typesetting small caps
% - \Lcase, \LCASE+, \LCASE- for typesetting case-sensitive forms
% - \Lfrac, \LFRAC+, \LFRAC- for typesetting inline fractions
% - \Lsalt, \LSALT+, \LSALT- for stylistic alternatives

% Customizable math alphabets for \mathcal, \mathfrak, \mathbb, and \mathscr
% (see https://ctan.org/pkg/mathalpha)
%% Relax some definitions already defined through loading
%% of libertinus-otf and re-defined by mathalpha
\let\mathbfcal\relax
\let\mathbfscr\relax
\let\mathbffrak\relax
\usepackage[
cal=stixtwoplain, % Calligraphic (\mathcal, \mathbfcal)
scr=stixtwofancy, % Script (\mathscr, \mathbfscr)
frak=stixtwo, % Fraktur (\mathfrak, \mathbffrak)
]{mathalpha}

% Micro-typographical enhancements
% (see https://ctan.org/pkg/microtype)
\usepackage[
tracking=smallcaps, % Tracking for small caps
verbose=errors, % All warnings become errors
]{microtype}

% Additional unicode symbols
\usepackage{pifont}

% Colors %
% Color facilities and extensions
% (see https://ctan.org/pkg/xcolor)
\usepackage[
luatex, % Use luatex driver
x11names, % Load predefined color set based on X11's colors
table, % Load tools for coloring table rows, columns, and cells
hyperref, % Integrate with hyperref
]{xcolor}

% Colorblind-friendly color schemes
% (see https://ctan.org/pkg/colorblind)
\usepackage[
Tol, % Load schemes by Paul Tol
pgf, % Enable colormaps for use with pgfplots (or TikZ)
]{colorblind}


% Boxes %
% Variable width minipages
% (see https://ctan.org/pkg/varwidth)
\usepackage{varwidth}

% Boxes with dashed lines
% (see https://ctan.org/pkg/dashbox)
\usepackage{dashbox}

% Stacking
% Customized stacking of objects
% (see https://ctan.org/pkg/stackengine)
\usepackage[
usestackEOL,% Use // for EOL symbol
]{stackengine}


% Figures %
% Handling of graphics, including both images and text
% (see https://ctan.org/pkg/graphicx)
\usepackage{graphicx}

% Specify directories to search for graphics files
\graphicspath{{figures/}}

% Programmatic creation of figures with TikZ
% (see https://tikz.dev/)
\usepackage{tikz}

% TikZ-based creation of trees
% (see https://www.ctan.org/pkg/forest)
\usepackage[%
edges,% Additional edge formatting specs
% external,% Use TikZ's externalization (see below), enhanced detection
]{forest}

% Cache ("externalize") TikZ figures in figures/tikz/
% to speed up compilation
% !! This requires the --shell-escape option to the engine !!
% (see https://tikz.dev/library-external)
\usetikzlibrary{external}
% \tikzexternalize[prefix=figures/tikz/]


% Tables %
% Enhanced functionality for array/tabular
% (see https://ctan.org/pkg/array)
\usepackage{array}

% Dynamic-width table columns
% (see https://ctan.org/pkg/tabularx)
\usepackage{tabularx}

% Multi-row cells, also adds functionality to makecell package (see below)
% (see https://ctan.org/pkg/multirow)
\usepackage{multirow}

% Column head layouts and multi-line cells
% (see https://ctan.org/pkg/makecell)
\usepackage{makecell}

% Support for page-crossing tables
% (see https://ctan.org/pkg/longtable)
\usepackage{longtable}

% Functionality for typesetting formal, typographically sound tables
% (see https://ctan.org/pkg/booktabs)
\usepackage{booktabs}

% Table environment separating style and content,
% providing amalgamation of features from other table
% packages, including the above
% (see https://ctan.org/pkg/tabularray)
\usepackage{tabularray}


% Mathematics %
% Tools for customizing theorem environments
% (see https://ctan.org/pkg/amsthm and https://ctan.org/pkg/thmtools)
\usepackage{amsthm}
\usepackage{thmtools}


% Algorithms and Code %
% Floating algorithm environment
% \usepackage{algorithm}

% Pseudocode typesetting (loads algorithmicx)
% (see https://ctan.org/pkg/algpseudocodex and https://ctan.org/pkg/algorithmicx)
\usepackage[
commentColor=T-Q-B0, % Comment color
rightComments=false, % Don't right justify comments
beginComment=\triangleright~, % Regular comment opening delimiter
beginLComment=\triangleright~, % Long comment opening delimiter
endLComment={}, % Long comment closing delimiter
]{algpseudocodex}

% Change typesetting of line numbers
\algrenewcommand{\alglinenumber}[1]{\textcolor{T-Q-B0}{\small{#1}:}}

% Define command that for comment with variable horizontal spacing
\algnewcommand{\Commentvs}[2][1em]{\hspace{#1}\Comment{#2}}

% Change some algorithmic keywords to shorter forms
\algrenewcommand{\algorithmicdo}{:}
\algrenewcommand{\algorithmicforall}{\textbf{forall}}
\algrenewcommand{\algorithmicprocedure}{\textbf{proc}}
\algrenewcommand{\algorithmicfunction}{\textbf{func}}

% Change algorithmic boxes to have dashed outline
\tikzset{algpxDefaultBox/.style={draw=dashed}}

% Patch algorithmic environment to prevent TikZ externalization
% \xpretocmd{\algorithmic}{\tikzexternaldisable}{}{\PatchFailed}%
% \xapptocmd{\endalgorithmic}{\tikzexternalenable}{}{\PatchFailed}%

% Algorithmic with tight top and bottom frame rule
\newsavebox{\algobox}
\newlength{\algorightmargin}
\NewDocumentEnvironment{algorithmicframed}{O{0} O{t} +b}%
{%
  \ifthenelse{\equal{#1}{0}}%
  {\setlength{\algorightmargin}{1em}}%
  {\setlength{\algorightmargin}{1.7em}}%
  % Save the body into a measuring box
  \begin{lrbox}{\algobox}
    \begin{varwidth}{\linewidth}
      \begin{algorithmic}[#1]
        #3
      \end{algorithmic}
    \end{varwidth}
  \end{lrbox}%
  % Typeset content with exact rule width
  \begin{minipage}[#2]{\wd\algobox + \algorightmargin}
    \rule[\ht\strutbox]{\linewidth}{0.4pt}\vspace{-\baselineskip}
    \begin{algorithmic}[#1]
    #3
    \end{algorithmic}
    \rule[\ht\strutbox]{\linewidth}{0.4pt}
  \end{minipage}%
}
{}


% Source code typesetting
% (see https://ctan.org/pkg/listings)
\usepackage{listings}

% Change title of listings list to match that for figures/tables
\renewcommand*{\lstlistlistingname}{List of Listings}

% Default listings setup
\lstset{%
  % Frame
  captionpos=t, % Position caption at top (mirroring what's typical for algorithms)
  frame=tb, % Top and bottom rules
  framesep=\smallskipamount, % Small skip between frame and listing content
  % Float placement
  floatplacement=tbhp, % Match float placement of built-in listing environment (as fallback)
  % Character printing and placement
  upquote=true, % Print backtick and single quote as is
  columns=[c]fixed, % Monospace characters, centered in their box
  keepspaces=true, % Don't drop spaces for column alignment
  tabsize=2, % Tabstops every 2 spaces
  mathescape=false, % Don't allow escaping to LaTeX with $
  showstringspaces=false, % Don't print characters for spaces
  % Line numbers
  numbers=none, % No line numbers
  % Basic style
  basicstyle={\normalsize\ttfamily},
}


% Citations %
\usepackage[
backend=biber, % Use biber as backend
bibencoding=utf-8, % Use UTF-8 encoding
bibstyle=alphabetic,citestyle=alphabetic, % Bibliography and citation styles
minnames=3,maxnames=4, % Allow up to 4 names in printed name lists, shorten to 3 if there are more
maxbibnames=1000, % Allow any number of names in bibliography entries
minitems=3,maxitems=4, % Allow up to 4 items in printed item lists, shorten to 3 if there are more
block=space, % Add some additional horizontal space between large parts of entry
hyperref=true, % Use hyperref for internal citation links
backref=true,backrefstyle=three+, % Add back-references in bibliography to pages with citations
labelalpha=true, % Request labelalpha and extraalpha data fields for customization
minalphanames=3,maxalphanames=4 % Allow up to 4 names in printed alpha citation labels, shorten to 3 if there are more
]{biblatex}

% Alphabetic-style citations should use lining (non-old) style numerals
\DeclareFieldFormat{labelalpha}{\liningnums{#1}}

% Force Oxford comma in lists in bibliography entries
\DeclareDelimFormat{finalnamedelim}{%
  \ifnumgreater{\value{liststop}}{2}{,}{}%
  \addspace\bibstring{and}\space}

% Make \fullcite macro write out all names (e.g., no "et al.")
\xpretocmd{\fullcite}{\AtNextCite{\defcounter{maxnames}{99}}}{}{\PatchFailed}


% Drafting %
% Comments (allows silencing of environments)
% (see https://ctan.org/pkg/comment)
\usepackage{comment}

% Dummy text
% (see https://ctan.org/pkg/lipsum)
\usepackage{lipsum}

% Todo notes
% (see https://ctan.org/pkg/todonotes)
\usepackage[
colorinlistoftodos=true,
color=T-Q-V4,
tickmarkheight=1ex,
textsize=small,
]{todonotes}

% Visualization of page layout
% (see https://ctan.org/pkg/layouts)
\usepackage{layouts}


% Order-dependent, late %
% Hyperlinking/referencing and friends
\usepackage[
pdfa=true, % Attempt compatibility with PDF/A
linktoc=page, % In TOC, have page number contain link (not text)
colorlinks=true, % Enable coloring of links (implicitly disabling boxes)
linkcolor={T-Q-MC6}, % Set link color to Tol's Medium Contrast Dark Red
citecolor={T-Q-MC6}, % Set cite color to Tol's Medium Contrast Dark Red
urlcolor={T-Q-MC6}, % Set url color to Tol's Medium Contrast Dark Red
]{hyperref}

% Cryptography-specific pseudocode typesetting
% (see https://ctan.org/pkg/cryptocode)
% No options are provided here, as they only seem to introduce
% macros for typesetting short cryptographic names or operators.
% We specify the ones we need ourselves in macrosetup.tex
\usepackage{cryptocode}

% Change default style (trying to synchronize with that of algorithmicx/algpseudocodex)
\pcsetargs{bodylinesep=0pt,codesize=\normalsize,subcodesize=\normalsize,jot=-\jot} % Distance between header line and bode, code size, and vertical spacing between lines

% Line numbers
\renewcommand*{\pclnstyle}[1]{\textcolor{T-Q-B0}{\textnormal{\small{#1}}}}
\renewcommand*{\pclnseparator}{\pclnstyle{:}}
\renewcommand*{\pclnspace}{0.5em}

% Comments
\renewcommand*{\pccomment}[2][1em]{\hspace{#1}\textcolor{T-Q-B0}{\triangleright~\textnormal{\textit{#2}}}}
\renewcommand*{\pclinecomment}[2][0em]{\pccomment[#1]{#2}}

% Keywords/constants/punctuation
\newcommand*{\pccolon}{%
  \ifmmode%
    \mathpunct{:}
  \else%
    :%
  \fi%
}
\renewcommand*{\pcelse}{\textnormal{\textbf{else}}}
\newcommand*{\pctern}[3]{%
  #1
  \mathbin{\textnormal{\textbf{if}}}
  #2
  \mathbin{\textnormal{\textbf{else}}}
  #3
}

% Clever referencing, e.g., automatic sorting, compression, context-awareness
% (see https://ctan.org/pkg/cleveref)
\usepackage[
capitalize, % Capitalize referenced names (e.g., Theorem instead of theorem)
noabbrev, % Don't use abbreviations of referenced names
]{cleveref}

% Force Oxford comma in reference lists
\newcommand*{\creflastconjunction}{, and\nobreakspace}
\newcommand*{\creflastgroupconjunction}{, and\nobreakspace}

% Fix cleveref use with cryptocode
\pcfixcleveref

%%% Local Variables:
%%% mode: LaTeX
%%% TeX-master: "../main.tex"
%%% TeX-engine: luatex
%%% End:
